\section{A useful answer before the company's history is complete}
\label{sec:real-package-connection}

The Apple examples now give us enough information to pose a more demanding
question: what can one connected answer say about the reported company and
its possible cash requirements? The answer must carry both the reported
amounts and the limits of the histories used to explain them. Otherwise a
complete-looking statement could conceal a missing liability, a hypothetical
retirement or an unavailable settlement date.

For the retained annual source, the selected statements report net income of
93,736 and other comprehensive income of 4,280, giving comprehensive income
of 98,016 million dollars. Assets of 364,980 less liabilities of 308,030 give
closing equity of 56,950. The following quarter reports net income of 36,330,
other comprehensive income of 383 and closing equity of 66,758
\citep[consolidated financial statements]{apple2024,apple2025q1}. These are
source-located readings of displayed amounts. Their units, periods, group
perimeter and statement populations accompany them. A share count or a
per-share amount keeps its own unit; it cannot enter a monetary total merely
because it sits next to one.

The connected result retains the 102 selected annual rows, 100 quarterly rows
and their separate opening balance-sheet comparisons. It also retains all
23 note dispositions, 70 known references, seven package obligations,
32 proposed source assertions and seven additional frontier entries from
the frozen inventory. A note's remaining facts stay attached to it. These
records preserve the declared package; they neither discover every possible
reference nor establish that every retained sentence has a formal meaning.

A query for the selected displayed statement tokens can therefore finish even
though the question about actual whole-company histories remains unresolved.
The family question specifies the issuer, sources and target histories, but
has no admitted member list or probability measure. The distinction matters:
an unresolved feasible remainder does not mean that no history exists, and
two illustrated histories do not mean that the real company has exactly two
possible histories. The combined answer is partial because the financial
explanation has obligations that the reading task does not.

Preserving a source token also requires restraint in arithmetic. The annual
cash-flow row for debt issuance displays a dash. Its observation remains a
reported dash. The retained financing reconciliation separately states the
narrow convention under which that dash contributes zero to its displayed
sum. This convention does not turn all dashes into zeros. Similarly, the
annual cash-flow opening amount of 30,737 differs from the opening balance-sheet
cash amount of 29,965 by 772. That population difference is an explanatory
residual, not permission to create a receipt or payment for 772.

\subsection{Matching annual rows leaves different liabilities to settle}

Two conditional asset histories illustrate why the connected result must
retain journals as well as totals. In H1, gross additions of 10,529 are paid
with 9,447 of cash and a capital payable of 1,082. Retirements remove gross
cost of 6,000 and accumulated depreciation of 5,636, leaving a loss of 364.
H2 adds and retires another 1,000 of gross cost: additions become 11,529,
retirements 7,000, loss 1,364 and capital payable 2,082. Depreciation is 8,200
in both branches. The extra acquisition and retirement cancel in net property,
plant and equipment; cash is unchanged. The asset journals alone leave H2
with 1,000 more liability and 1,000 less profit and equity.

Now suppose a distinct non-tax expense liability exists: its amount is $q$
in H1 and $q-1{,}000$ in H2, with $q\geq1{,}000$. Each branch debits expense
and credits that liability. This is an independent economic premise, not an
event inferred from a statement residual. Suppose also that the retirement
loss and this expense share the product-cost row, both liabilities share the
same other-current-liability bucket, the extra asset cost has the same gross
class, and current and deferred tax, all other activities, shares and
distributions are unchanged. Then the expense difference of $-1{,}000$
offsets the extra retirement loss, and the liability difference offsets the
extra capital payable.

The indirect cash-flow presentation still needs a specified allocation. Let
$\alpha\in[0,1]$ allocate the expense-liability adjustment once between the
noncash adjustment and the working-capital adjustment. If $\Delta$ means
H2 minus H1, the two displayed-row differences become
\begin{align}
\Delta\text{other adjustments}
 &=1{,}000+\alpha(-1{,}000)=1{,}000(1-\alpha),\\
\Delta\text{other liabilities}
 &=(1-\alpha)(-1{,}000)=-1{,}000(1-\alpha).
\label{eq:real-package-cash-presentation}
\end{align}
The two differences sum to zero for every allowed $\alpha$. At $\alpha=1$,
each is zero and all 102 selected annual row differences vanish under these
assumptions. At $\alpha=0$ or $1/2$, the same total operating cash has a
different presentation in those two rows. Here $\alpha$ is a reporting
assumption, not a probability over histories. Matching differences also does
not construct an absolute feasible remainder for the company. The illustrative
journals remain conditional subsystem histories.

Their future consequences remain joint. With $q=2{,}000$, carrying-value
settlement of the two liabilities gives the following changes, in millions,
if both payments occur in the first future period:
\begin{center}
\begin{tabular}{lrrr}
\toprule Branch & Operating cash & Investing cash & Total cash\\\midrule
H1 & $-2{,}000$ & $-1{,}082$ & $-3{,}082$\\
H2 & $-1{,}000$ & $-2{,}082$ & $-3{,}082$\\\bottomrule
\end{tabular}
\end{center}
Paying the capital liability is an investing movement; paying the separately
accrued expense liability is an operating movement. Each payment debits its
liability and credits cash. Under the explicit assumption of no incremental
tax, currency or interest effects, it creates no new profit, other
comprehensive income or equity movement. Sufficient unrestricted funding is
assumed; no opening cash balance is manufactured.

Taking the lowest operating amount from H1 and the lowest investing amount
from H2 would produce $-4{,}082$. Neither history produces that total. The
result must return the two pairs, not independently selectable marginal
endpoints. Moving capital payment to the second period changes first-period
cash to $-2{,}000$ in H1 and $-1{,}000$ in H2; moving the expense payment
instead changes it to $-1{,}082$ and $-2{,}082$. Thus equal eventual cash does
not identify its timing. The six specified paths retain both branch identity
and payment period.

When dates are missing, the response keeps the conditional eventual amounts
and marks period timing unresolved. It distinguishes these calculated amounts
from reported tokens. A request for future depreciation remains unsupported
without the asset cohorts, tax state and other required facts. A probability
request remains unsupported without an admitted law. Neither refusal erases
the amounts that a narrower question can answer.

\subsection{The information used belongs to a particular view}

The later annual source provides a different kind of progress. Its FY2024
Americas comparative shows sales of 167,045, cost of sales of $-89,587$ and
selling and marketing expense of $-9,802$, giving operating income of 67,656
\citep[Note 13, FY2024 Americas comparative]{apple2025annual}.
The selected original source lacks those two separate expense rows. An
original-source query therefore returns unresolved components, while a query
explicitly naming the later source can return them. The result links the later
view separately; it does not silently enlarge the original view or assert
that measurement scope was unchanged. Later detail is not, by itself, an
accounting-error correction.

Named-source access answers what retained bytes contain. Investor-information
and system-as-of queries additionally require accepted clocks for those exact
bytes. Where publication, ingestion or review time is unestablished, the
historical query remains partial with a named timing obligation. A filing
event alone cannot date a separately converted PDF byte for byte. This prevents
the later comparative detail from becoming evidence that an earlier investor
could already have used.

JPMorgan's reference conflict illustrates the same discipline in another
company. In the retained 2023 report, Item 15 faithfully says page 156; Item 8
and the auditor heading support printed page 163, physical page 165, for
navigation \citep[Items 8 and 15; printed pp.~156 and 163]{jpmorgan2023tenk}.
The wrong target is printed page 156, physical page 158. Both use
the same two-page offset, so changing the offset cannot reconcile the targets.
A proposed navigation correction preserves the literal statement. If a
counterfactual earlier view had accepted the wrong target, its selected page
set, bounded search and coverage summary would need reconsideration, while
the literal and an independently sourced statement amount could remain usable.
A one-page negative search could never establish absence from the whole filing.

That review proposal has its own issuer and source dependencies. It supplies
no Apple financial constraint, no posting permission and no observed human
decision. The connected financial result and the separate review proposal can
therefore be inspected together without joining their companies' histories.
Section~\ref{sec:real-package-reference} gives the exact exchange locations
and rejection boundaries that preserve these distinctions.
